\section{Which questions prices answer cheaply}\label{app:questions}

The cost results of Section~\ref{sec:cost} concern composition. This section asks the same question of every linear question about the population, and finds a dichotomy: some questions are answered at the usual statistical rate whatever the number of groups and the spacing of their memories, some cannot be answered at all without a bound on memory, and the cost of each is the value of a problem of polynomial approximation. The results and proofs are adapted from the first author's thesis (in preparation).

\paragraph{Setting.} Write $M=\sum_iw_i\delta_{r_i}$ for the merged measure, so that $k_m=\int r^m\,dM$, and $\mathcal M(\bar r,\bar W)$ for the positive measures on $[0,\bar r]$ of total mass at most $\bar W$, with $\bar r<1$ a memory bound. A \emph{linear question} is $F_g(M)=\int g\,dM$ with $g$ continuous on $[0,\bar r]$; examples are a single tap ($g(r)=r^m$), the discounted feedback ($g(r)=(1-\beta r)^{-1}$, $\beta<1$), and the total feedback weight $\Phi=\int dM/(1-r)=\sum_m k_m$, which with $a$ and $\kappa$ fixes the feedback budget $S=a+\kappa\Phi$. The observer runs $N$ independent experiments, each returning the first $L$ kernel taps with independent $N(0,\sigma^2)$ noise. One way to obtain them is a steering probe that holds the price at a unit impulse (as in the proof of Theorem~\ref{thm:anypop}), with aggregate feedback flow observed: the flow at date $t+1$ is then $k_t$. Price records are discussed at the end of the section. For accuracy $\epsilon>0$ and level $1-\alpha$, $N_{\rm cert}(\epsilon;F_g)$ is the least $N$ for which some interval built from the data has half-width at most $\epsilon$ and covers $F_g(M)$ with probability at least $1-\alpha$ for every $M\in\mathcal M(\bar r,\bar W)$.

For $g(r)=\sum_ta_tr^t$ write $\tau_g(T)=\bar W\sup_{r\le\bar r}|\sum_{t\ge T}a_tr^t|$, $A_g(T)=\sum_{t<T}a_t^2$, $T_g(\epsilon)=\min\{T\ge1:\tau_g(T)\le\epsilon/2\}$, and $E_T(g)=\min_{\deg p<T}\max_{[0,\bar r]}|g-p|$. For $c\in\mathbb R^L$ put $p_c(r)=\sum_{t<L}c_tr^t$, and for a function $h$ on $[0,\bar r]$ put $\mathcal R(h)=(\sup h)_++(-\inf h)_+$, which is at least $\|h\|_\infty$. The modulus of the question is
\[
 \omega_L(\delta)=\sup\bigl\{F_g(M_1)-F_g(M_0):\ M_0,M_1\in\mathcal M(\bar r,\bar W),\ \|k_{0:L}(M_1)-k_{0:L}(M_0)\|_2\le\delta\bigr\},
\]
the largest difference in the answer between two populations whose observed taps differ by at most $\delta$.

\begin{theorem}[The cost of a linear question]\label{thm:questions}
Let $z_\beta$ be the standard normal $\beta$-quantile and $\alpha\in(0,1/2)$.
\begin{enumerate}[label=(\roman*)]
\item \emph{Upper bound.} If $L\ge T_g(\epsilon)$, then $N_{\rm cert}(\epsilon;F_g)\le\lceil4z_{1-\alpha/2}^2\sigma^2A_g(T_g(\epsilon))\epsilon^{-2}\rceil$.
\item \emph{Truncation is necessary.} For every $T$ there are $M,M'\in\mathcal M(\bar r,\bar W)$ whose first $T$ taps coincide and with $|F_g(M)-F_g(M')|\ge2\bar WE_T(g)$; any certificate that uses only the first $T$ taps has worst-case half-width at least $\bar WE_T(g)$.
\item \emph{Long memory makes some questions impossible.} On the positive measures on $[0,1)$ of mass at most $\bar W$, with no memory bound, if $\limsup_{r\uparrow1}|g(r)|\sqrt{1-r}=\infty$, then for $\alpha<1/4$ no finite $N$ certifies $F_g$ to any finite accuracy.
\item \emph{Exact cost.} For every depth $L$ and $\delta>0$,
\[
 \omega_L(\delta)=\min_{c\in\mathbb R^L}\bigl\{\delta\|c\|_2+\bar W\,\mathcal R(g-p_c)\bigr\},
\]
and $\omega_L(0)=2\bar WE_L(g)$. If $N$ experiments certify $F_g$ to half-width $\epsilon$ at level $1-\alpha$, then $\omega_L(2z_{1-\alpha}\sigma/\sqrt N)\le2\epsilon$; conversely, if $\omega_L(2z_{1-\alpha/2}\sigma/\sqrt N)\le2\epsilon$, an affine interval $\sum_tc_t\bar k_t+b\pm\epsilon$ certifies $F_g$ with $N$ experiments. With $N_\omega(\epsilon)=\min\{N:\omega_L(2z_{1-\alpha}\sigma/\sqrt N)\le2\epsilon\}$ and $\lambda_\alpha=z_{1-\alpha/2}/z_{1-\alpha}$,
\[
 N_\omega(\epsilon)\le N_{\rm cert}(\epsilon;F_g)\le N_\omega(\epsilon/\lambda_\alpha).
\]
\end{enumerate}
\end{theorem}

\begin{proof}
(i) Let $\bar k_t$ be the average of the $N$ replicates of tap $t$, and $T=T_g(\epsilon)$. The estimator $\hat F=\sum_{t<T}a_t\bar k_t$ has mean $\sum_{t<T}a_tk_t$ and variance $\sigma^2A_g(T)/N$. Its bias is $|\sum_{t\ge T}a_tk_t|=|\int\sum_{t\ge T}a_tr^t\,dM|\le\tau_g(T)\le\epsilon/2$ for every $M$ in the class. So $\hat F\pm(\epsilon/2+z_{1-\alpha/2}\sigma\sqrt{A_g(T)/N})$ covers $F_g(M)$ with probability at least $1-\alpha$, and its half-width is at most $\epsilon$ once $z_{1-\alpha/2}\sigma\sqrt{A_g(T)/N}\le\epsilon/2$.

(ii) By duality for best uniform approximation, $E_T(g)=\max\int g\,d\nu$ over signed measures $\nu$ on $[0,\bar r]$ of total variation at most one that annihilate every polynomial of degree below $T$. Take a maximizer and write $\nu=\nu^+-\nu^-$. Orthogonality to constants gives $\nu^+([0,\bar r])=\nu^-([0,\bar r])=:c\le1/2$. The measures $M=(\bar W/c)\nu^+$ and $M'=(\bar W/c)\nu^-$ lie in the class, their first $T$ moments, which are their first $T$ taps, coincide, and $F_g(M)-F_g(M')=(\bar W/c)E_T(g)\ge2\bar WE_T(g)$. Two populations that produce identically distributed data cannot both be covered by an interval of half-width below half the difference of their answers.

(iii) Compare $M_0=0$ with $M_1=c\,\delta_r$. Under $M_1$ the taps are $k_t=cr^t$, so the Kullback--Leibler divergence between the two data laws over any number of taps is at most $Nc^2/(2\sigma^2(1-r^2))$. Take $c=\sigma\sqrt{(1-r^2)/N}$, which is at most $\bar W$ for $r$ near one; the divergence is at most $1/2$, so by Pinsker's inequality the total variation distance is at most $1/2$. The answers differ by $c|g(r)|\ge\sigma\sqrt{(1-r)/N}\,|g(r)|$, which is unbounded along a sequence $r\uparrow1$. If an interval of half-width $\epsilon$ covered $F_g$ with probability at least $1-\alpha$ under both, then once the answers differ by more than $2\epsilon$ the event that the interval contains $F_g(M_0)$ has probability at least $1-\alpha$ under $M_0$ and at most $\alpha$ under $M_1$, so the total variation distance would be at least $1-2\alpha>1/2$.

(iv) \emph{Duality.} For any $c\in\mathbb R^L$ and any admissible pair, with $h=g-p_c$,
\[
 F_g(M_1)-F_g(M_0)=\int h\,dM_1-\int h\,dM_0+\langle c,k_{0:L}(M_1)-k_{0:L}(M_0)\rangle\le\bar W\mathcal R(h)+\delta\|c\|_2,
\]
because $\int h\,dM_1\le\bar W(\sup h)_+$ and $-\int h\,dM_0\le\bar W(-\inf h)_+$ for measures of mass at most $\bar W$. So $\omega_L(\delta)$ is at most the minimum. For the reverse inequality let $\mu=(M_0,M_1)$ and $\Lambda(\mu,c)=\int h\,d(M_1-M_0)+\delta\|c\|_2$. Since $\inf_c\{\delta\|c\|_2-\langle c,v\rangle\}$ is $0$ if $\|v\|_2\le\delta$ and $-\infty$ otherwise, $\omega_L(\delta)=\sup_\mu\inf_c\Lambda(\mu,c)$. The set $\mathcal M(\bar r,\bar W)^2$ is convex and weak-$*$ compact, $\Lambda$ is affine and weak-$*$ continuous in $\mu$ (as $h$ is continuous) and convex and continuous in $c$, so Sion's minimax theorem gives $\omega_L(\delta)=\inf_c\sup_\mu\Lambda(\mu,c)$. The inner supremum is $\bar W\mathcal R(g-p_c)+\delta\|c\|_2$, attained by mass $\bar W$ at a maximizer of $h$ if $\sup h>0$ and at a minimizer if $\inf h<0$. For $\delta>0$ the objective is convex, continuous and coercive in $c$, so the infimum is attained. At $\delta=0$ the same argument gives $\omega_L(0)=\inf_c\bar W\mathcal R(g-p_c)$; adding a constant to $p_c$ makes $g-p_c$ change sign, where $\mathcal R$ is the oscillation, so the infimum is $\bar W$ times the least oscillation of $g-p$ over polynomials of degree below $L$, which is $2\bar WE_L(g)$.

\emph{Lower bound.} The averages $\bar k=(\bar k_0,\ldots,\bar k_{L-1})$ are sufficient and distributed as $N(k_{0:L}(M),(\sigma^2/N)I_L)$, so the total variation distance between the data laws of two populations is $2\Phi_{\mathcal N}(\sqrt N\|\Delta k\|_2/(2\sigma))-1$, with $\Phi_{\mathcal N}$ the standard normal distribution function. For an admissible pair with $\|\Delta k\|_2<2z_{1-\alpha}\sigma/\sqrt N$ this is below $1-2\alpha$. If their answers differed by more than $2\epsilon$, an interval of half-width $\epsilon$ covering each with probability $1-\alpha$ would separate the two laws with total variation at least $1-2\alpha$, as in (iii). Hence $\omega_L(t)\le2\epsilon$ for every $t<2z_{1-\alpha}\sigma/\sqrt N$. The modulus is concave (the admissible pairs form a convex set and the constraint is a norm), finite and nondecreasing, hence continuous on $(0,\infty)$, so $\omega_L(2z_{1-\alpha}\sigma/\sqrt N)\le2\epsilon$.

\emph{Upper bound.} Put $\delta=2z_{1-\alpha/2}\sigma/\sqrt N$, let $c$ attain the minimum in the dual form and $h=g-p_c$. The interval $\sum_tc_t\bar k_t+b\pm\epsilon$ with $b=\bar W[(\sup h)_+-(-\inf h)_+]/2$ has bias at most $\bar W\mathcal R(h)/2$ for every population, by the bounds in the duality step, and normal noise with standard deviation $\sigma\|c\|_2/\sqrt N$. It covers $F_g$ with probability at least $1-\alpha$ whenever $\epsilon\ge\bar W\mathcal R(h)/2+z_{1-\alpha/2}\sigma\|c\|_2/\sqrt N=\omega_L(\delta)/2$.

\emph{Sandwich.} The left inequality restates the lower bound. Concavity and $\omega_L(0)\ge0$ give $\omega_L(\lambda t)\le\lambda\omega_L(t)$ for $\lambda\ge1$. If $\omega_L(2z_{1-\alpha}\sigma/\sqrt N)\le2\epsilon/\lambda_\alpha$, then $\omega_L(2z_{1-\alpha/2}\sigma/\sqrt N)\le2\epsilon$, and the upper bound applies.
\end{proof}

Part~(iv) has the structure of the classical theory of estimating a linear functional over a convex class \citep{donoho1994}: one modulus governs both the lower bound and an affine interval. What is specific to this market is the dual form. The cost of a question about the traders is the value of a problem of polynomial approximation of $g$ on the memory range, in which the coefficients are penalized because each multiplies a noisy tap. The number of groups and the spacing of their memories do not enter.

\begin{proposition}[Which questions are parametric]\label{prop:parametric}
Let $\omega_\infty(\delta)=\inf_L\omega_L(\delta)$. Then $\sup_{0<\delta\le1}\omega_\infty(\delta)/\delta<\infty$, so that $N_{\rm cert}(\epsilon;F_g)=O(\sigma^2\epsilon^{-2})$ with observation depth $L$ of order $\log(1/\epsilon)$, if and only if $g$ agrees on $[0,\bar r]$ with a power series $\sum_ta_tr^t$ whose coefficients are square summable. When $g$ is a power series on $[0,\bar r]$ with $\bar r>0$, this is the condition $\sum_ta_t^2<\infty$ on its Taylor coefficients.
\end{proposition}

\begin{proof}
If $a\in\ell^2$, take $c=(a_0,\ldots,a_{L-1})$. By Cauchy--Schwarz, $\mathcal R(g-p_c)\le2\sup_{r\le\bar r}|\sum_{t\ge L}a_tr^t|\le2\|a\|_2\bar r^L/\sqrt{1-\bar r^2}$, so $\omega_L(\delta)\le\delta\|a\|_2+2\bar W\|a\|_2\bar r^L/\sqrt{1-\bar r^2}$, which is at most $C\delta$ once $L\ge\log(1/\delta)/\log(1/\bar r)$; Theorem~\ref{thm:questions}(iv) then gives the parametric cost. Conversely, suppose $\omega_\infty(\delta)\le C\delta$ for $\delta\in(0,1]$. For $\delta_n=1/n$ choose $L_n$ and $c^{(n)}\in\mathbb R^{L_n}$ with $\delta_n\|c^{(n)}\|_2+\bar W\mathcal R(g-p_{c^{(n)}})\le2C\delta_n$. Then $\|c^{(n)}\|_2\le2C$ and $\|g-p_{c^{(n)}}\|_\infty\le2C\delta_n/\bar W\to0$. Padded with zeros, the $c^{(n)}$ are bounded in $\ell^2$, so a subsequence converges weakly to some $a\in\ell^2$. For each $r\in[0,\bar r]$ the sequence $(r^t)_t$ lies in $\ell^2$, so $p_{c^{(n)}}(r)=\langle c^{(n)},(r^t)\rangle\to\sum_ta_tr^t$ along the subsequence, while $p_{c^{(n)}}(r)\to g(r)$. Power series that agree on an interval are equal, which gives the last sentence.
\end{proof}

\paragraph{Where the questions of this paper sit.} A single tap and the discounted feedback have square-summable coefficients and are parametric: they cost of order $\sigma^2\epsilon^{-2}$ experiments however many groups there are and however close their memories. Square-summable power series are the restrictions of the Hardy space $H^2$ of the unit disc, whose functions grow at most like $(1-r)^{-1/2}$ as $r\uparrow1$; that is exactly the growth beyond which part~(iii) rules out certification without a memory bound. The total feedback weight $\Phi$, which fixes the feedback budget and with it stability and worst-case transient loss, has $a_t\equiv1$ and sits on the boundary: \citet{niroula2026b} show that it costs $\sigma^2\epsilon^{-2}\log(1/\epsilon)$ experiments and that the logarithm cannot be removed. Composition sits at the far end. The weight of a group is not a linear functional of $M$ at all, and by Theorem~\ref{thm:kcohort} its cost grows like $(\Delta\eta)^{-(4K-2)}$. Market-wide questions are cheap; questions about who trades are not.

\paragraph{Price records.} The theorem is stated for observed kernel taps. With price records only, the first $T+1$ price taps are a function of the first $T$ kernel taps ($h_0=\kappa$ and $h_{t+1}=ah_t+\kappa\sum_{j\le t}k_jh_{t-j}$), so part~(ii) holds with $T+1$ price taps in place of $T$ kernel taps, and every lower bound for kernel taps that rests on equal taps holds for price taps. In the other direction, the map from the first $L+1$ price taps to the first $L$ kernel taps is triangular, and its linearization at a stable economy is a lower-triangular Toeplitz operator whose symbol, the reciprocal of the square of the closed-loop transfer function up to a factor, is bounded on the unit circle uniformly in $L$. Estimated kernel taps therefore have $\ell^2$ errors of order $\sigma N^{-1/2}$ to first order, and the parametric case of Proposition~\ref{prop:parametric} carries over to price records to first order, on classes with a uniform stability margin. We do not state the higher-order terms.

\paragraph{Numerical check.} With two groups at memories $0.5\pm\Delta\eta$, the error of the truncated estimator of the discounted feedback ($\beta=0.9$) from noisy taps does not depend on $\Delta\eta$, while the error of the weights by Prony's method grows like $(\Delta\eta)^{-3}$: fitted slopes $-0.02$ and $-3.02$ over $\Delta\eta\in\{0.2,0.1,0.05\}$, 2,000 replications each (\nolinkurl{research_extension/theorem_checks/mc_questions.py}).
